% GENERATED FILE. Edit canonical lessons, then run npm run generate:books.
% canonical-source-hash: fnv1a64:6bc99960d8028f68
% canonical-lessons: FR-C278-actualite, FR-C278-presse, FR-C278-quotidien, FR-C278-hebdomadaire, FR-C278-abonne

\chapter{La presse}
\label{ch:fr-la-presse}
\hlchaptermodality{278}

\begin{chapteropening}
\textbf{By the end of this chapter:} \emph{I can talk about the news and the papers: current events, the press, daily and weekly papers, and their subscribers.}

Complete the last lesson of chapter 278: I can talk about the news and the papers: current events, the press, daily and weekly papers, and their subscribers.
\end{chapteropening}

\section[l'actualité]{l'actualité — the news, current events}

\label{lesson:FR-C278-actualite}



\begin{quote}
Before the new one: say the French for the instant, then the French for at the end of.
\end{quote}

\subsection*{You'll want to know: l'actualité}
\textbf{l'actualité} — \textquotedblleft{}the news, current events\textquotedblright{}.

\begin{quote}
Mon père suit l'actualité tous les matins.
\end{quote}

\begin{quote}
\emph{My father follows the news every morning.}
\end{quote}

\emph{Les actualités} — the news on TV or radio, a little old-fashioned. \emph{D'actualité} — topical.

\begin{cousinweb}
\emph{actualité} is a learned borrowing from Medieval Latin \emph{āctuālitās} (EtymDB), built on \emph{āctuālis} \textquotedblleft{}active, real\textquotedblright{}, from \emph{agere} \textquotedblleft{}to do\textquotedblright{} (Lewis \& Short): what is happening now.
\end{cousinweb}

\begin{grammarlens}[title={What you've built}]
One more word for telling, explaining and describing at B1.
\end{grammarlens}

\subsection*{Your turn}
\begin{itemize}
\raggedright
  \item \emph{Say it:} \emph{l'actualité}
  \item \emph{Say it:} \emph{l'actualité}, once more
  \item \emph{Say it:} a sentence of your own with \emph{l'actualité}
  \item \emph{Recall:} say the French for the instant, then the French for at the end of, then say \emph{l'actualité} again
\end{itemize}

\begin{tcolorbox}[breakable,colback=teal!4,colframe=teal!35!black,title={Before you move on}]
What does l'actualité mean? (\textquotedblleft{}The news, current events\textquotedblright{}.) Say it once more.
\end{tcolorbox}
\section[la presse]{la presse — the press, the newspapers}

\label{lesson:FR-C278-presse}



\begin{quote}
Before the new one: say the French for at the end of, then the French for the news.
\end{quote}

\subsection*{You'll want to know: la presse}
\textbf{la presse} — \textquotedblleft{}the press, the newspapers\textquotedblright{}.

\begin{quote}
La presse a beaucoup parlé de cette histoire.
\end{quote}

\begin{quote}
\emph{The press talked a lot about this story.}
\end{quote}

\emph{La presse écrite} — the print press. It also names a machine that presses.

\begin{cousinweb}
\emph{presse} is built on \emph{presser} \textquotedblleft{}to press\textquotedblright{}, which is inherited from Latin \emph{pressāre}, built on \emph{premere} \textquotedblleft{}to press\textquotedblright{} (Lewis \& Short): first the printing press, then the papers it printed.
\end{cousinweb}

\begin{grammarlens}[title={What you've built}]
One more word for telling, explaining and describing at B1.
\end{grammarlens}

\subsection*{Your turn}
\begin{itemize}
\raggedright
  \item \emph{Say it:} \emph{la presse}
  \item \emph{Say it:} \emph{la presse}, once more
  \item \emph{Say it:} a sentence of your own with \emph{la presse}
  \item \emph{Recall:} say the French for at the end of, then the French for the news, then say \emph{la presse} again
\end{itemize}

\begin{tcolorbox}[breakable,colback=teal!4,colframe=teal!35!black,title={Before you move on}]
What does la presse mean? (\textquotedblleft{}The press, the newspapers\textquotedblright{}.) Say it once more.
\end{tcolorbox}
\section[le quotidien]{le quotidien — the daily paper; daily life}

\label{lesson:FR-C278-quotidien}



\begin{quote}
Before the new one: say the French for the news, then the French for the press.
\end{quote}

\subsection*{You'll want to know: le quotidien}
\textbf{le quotidien} — \textquotedblleft{}the daily paper; daily life\textquotedblright{}.

\begin{quote}
Il achète un quotidien tous les matins.
\end{quote}

\begin{quote}
\emph{He buys a daily paper every morning.}
\end{quote}

As an adjective, \emph{quotidien} means daily: \emph{la vie quotidienne}.

\begin{cousinweb}
\emph{quotidien} is a learned borrowing from Latin \emph{quotīdiānus}, \emph{cottīdiānus} \textquotedblleft{}daily\textquotedblright{} (EtymDB), from \emph{cottīdiē} \textquotedblleft{}every day\textquotedblright{}, built on \emph{diēs} \textquotedblleft{}day\textquotedblright{} (Lewis \& Short).
\end{cousinweb}

\begin{grammarlens}[title={What you've built}]
One more word for telling, explaining and describing at B1.
\end{grammarlens}

\subsection*{Your turn}
\begin{itemize}
\raggedright
  \item \emph{Say it:} \emph{le quotidien}
  \item \emph{Say it:} \emph{le quotidien}, once more
  \item \emph{Say it:} a sentence of your own with \emph{le quotidien}
  \item \emph{Recall:} say the French for the news, then the French for the press, then say \emph{le quotidien} again
\end{itemize}

\begin{tcolorbox}[breakable,colback=teal!4,colframe=teal!35!black,title={Before you move on}]
What does le quotidien mean? (\textquotedblleft{}The daily paper; daily life\textquotedblright{}.) Say it once more.
\end{tcolorbox}
\section[l'hebdomadaire]{l'hebdomadaire — the weekly (paper or magazine)}

\label{lesson:FR-C278-hebdomadaire}



\begin{quote}
Before the new one: say the French for the press, then the French for the daily paper.
\end{quote}

\subsection*{You'll want to know: l'hebdomadaire}
\textbf{l'hebdomadaire} — \textquotedblleft{}the weekly (paper or magazine)\textquotedblright{}.

\begin{quote}
Ma grand-mère lit un hebdomadaire sur la cuisine.
\end{quote}

\begin{quote}
\emph{My grandmother reads a weekly about cooking.}
\end{quote}

Often \emph{un hebdo} in speech. As an adjective: \emph{une réunion hebdomadaire}.

\begin{cousinweb}
\emph{hebdomadaire} is a learned borrowing from Late Latin \emph{hebdomadārius} (EtymDB), built on Greek \emph{hebdomás} \textquotedblleft{}a group of seven, a week\textquotedblright{}.
\end{cousinweb}

\begin{grammarlens}[title={What you've built}]
One more word for telling, explaining and describing at B1.
\end{grammarlens}

\subsection*{Your turn}
\begin{itemize}
\raggedright
  \item \emph{Say it:} \emph{l'hebdomadaire}
  \item \emph{Say it:} \emph{l'hebdomadaire}, once more
  \item \emph{Say it:} a sentence of your own with \emph{l'hebdomadaire}
  \item \emph{Recall:} say the French for the press, then the French for the daily paper, then say \emph{l'hebdomadaire} again
\end{itemize}

\begin{tcolorbox}[breakable,colback=teal!4,colframe=teal!35!black,title={Before you move on}]
What does l'hebdomadaire mean? (\textquotedblleft{}The weekly (paper or magazine)\textquotedblright{}.) Say it once more.
\end{tcolorbox}
\section[l'abonné]{l'abonné — the subscriber}

\label{lesson:FR-C278-abonne}



\begin{quote}
Before the new one: say the French for the daily paper, then the French for the weekly.
\end{quote}

\subsection*{You'll want to know: l'abonné}
\textbf{l'abonné} — \textquotedblleft{}the subscriber\textquotedblright{}.

\begin{quote}
Les abonnés reçoivent le journal à la maison.
\end{quote}

\begin{quote}
\emph{Subscribers get the paper at home.}
\end{quote}

\emph{Être abonné à un journal} — to take a paper. Feminine: \emph{l'abonnée}. Online: \emph{les abonnés} are followers.

\begin{cousinweb}
\emph{abonné} is the past participle of \emph{abonner} \textquotedblleft{}to subscribe\textquotedblright{} (EtymDB), built on Old French \emph{bonne}, \emph{borne} \textquotedblleft{}a boundary\textquotedblright{}, which is usually traced to a Gaulish word: a price fixed within limits.
\end{cousinweb}

\begin{grammarlens}[title={What you've built}]
One more word for telling, explaining and describing at B1.
\end{grammarlens}

\subsection*{Your turn}
\begin{itemize}
\raggedright
  \item \emph{Say it:} \emph{l'abonné}
  \item \emph{Say it:} \emph{l'abonné}, once more
  \item \emph{Say it:} a sentence of your own with \emph{l'abonné}
  \item \emph{Recall:} say the French for the daily paper, then the French for the weekly, then say \emph{l'abonné} again
\end{itemize}

\begin{tcolorbox}[breakable,colback=teal!4,colframe=teal!35!black,title={Before you move on}]
What does l'abonné mean? (\textquotedblleft{}The subscriber\textquotedblright{}.) Say it once more.
\end{tcolorbox}
